%% kapa: exam=1 round=27 year=2016 period=2 form=A subject=appraisal_law subject_ko=감정평가관계법규 seq=26 number=26
%% source: Q-Net 감정평가사 기출문제 (cst00309) artlSeq=5169402 file=extracted/1008/5169402/제27회 감정평가사 1차 문제/제27회 감정평가사 2교시문제(A형).pdf
%% conversion: kordoc 4.24.51 -> tools/gen1.py (auto); review=proofread-sonnet; flags=none
\begin{question}{26}
건축법령상 위반 건축물 등에 대한 조치에 관한 설명으로 옳지 않은 것은?
\begin{choices}
  \item 허가권자는 건축물이 건축법령에 위반되는 경우 그 건축물의 현장관리인에게 공사의 중지를 명할 수 있다.
  \item 건축물이 용적률을 초과하여 건축된 경우 해당 건축물에 적용되는 시가표준액의 100분의 10에 해당하는 금액으로 이행강제금이 부과된다.
  \item 허가권자는 이행강제금을 부과하기 전에 이행강제금을 부과․징수한다는 뜻을 미리 문서로써 계고(戒告)하여야 한다.
  \item 허가권자는 이행강제금 부과처분을 받은 자가 이행강제금을 납부기한까지 내지 아니하면「지방세외수입금의 징수 등에 관한 법률」에 따라 징수한다.
  \item 허가권자는 시정명령을 받은 자가 이를 이행하면 새로운 이행강제금의 부과를 즉시 중지하되, 이미 부과된 이행강제금은 징수하여야 한다.
\end{choices}
\end{question}
\answer{2}
